\subsection{Overview}

Across the literature on agricultural subsurface (tile) drainage, authors converge on a consistent message: while individual practices and processes are increasingly well characterised, the field still lacks the long-term, multi-scale, and economically grounded evidence needed to design, target, and deploy cost-effective nutrient-mitigation strategies at the landscape scale. This section synthesises the research frontier around several recurring lines of inquiry. First, there is a widespread call for longer-duration, multi-year, and multi-site monitoring to smooth climatic variability and capture the true magnitude of system responses \parencite{Randall2003,Randall2005,Waring2024,Jeong2018,Cooke2012}. Second, nitrogen management remains a primary focus for refinement of application rates, timing, and decision-support tools \parencite{Christianson2015b,Bandura2025,Liu2021}. Third, phosphorus transport in subsurface drainage is repeatedly flagged as under-studied relative to its importance, particularly for dissolved and bioavailable fractions \parencite{Hertzberger2019,Yuan2022,deOliveira2022}. Fourth, a strong theme is that no single best management practice (BMP) is a "silver bullet": bundling in-field, edge-of-field, and water-management practices, and targeting them to hydrologically active critical source areas, is essential \parencite{Sunohara2016,Lenhart2016,Christianson2014,Baulch2019}. Finally, authors emphasise that biophysical understanding must be coupled with economics, incentive design, policy reform, modelling advances, and stakeholder engagement before mitigation can be scaled \parencite{Liu2021,Kvítek2023,David2015}. The subsections below organise these open questions and priorities in turn, retaining the specific quantitative targets, methods, and conditions proposed by the source studies.

\subsection{Long-term, multi-scale monitoring and experimental design}

A dominant frontier priority is the need for longer observation periods to disentangle management effects from climatic noise. \textcite{Randall2003,Randall2005} argue that long-term studies with detailed measurements are necessary to smooth year-to-year climatic variability and thereby yield more meaningful information, and stress that both the corn and soybean phases of a rotation should be evaluated when assessing nitrogen (N) management effects on nitrate losses to subsurface drainage. \textcite{Waring2024} similarly conclude that long-term studies are essential to capture a system's response to conservation practices and the true magnitude of that response, emphasising site-specific impacts. \textcite{Cooke2012} note that studies of longer duration, and many more years of yield data, are required to determine the effect of drainage water management (DWM) on yield while eliminating confounding factors. \textcite{Pittelkow2017} likewise highlight the need to evaluate N losses over multiple years on an annualised basis. Continuous and comprehensive monitoring of N application rate, timing, and method on nitrate loss and crop yield in tile-drained landscapes is advocated to give modelling a more rational setup \parencite{Jeong2018}.

Several studies extend this call to non-growing-season and multi-year hydrologic coverage. \textcite{VanEsbroeck2016} recommend collecting field data during the non-growing season when estimating phosphorus (P) loads and assessing BMP efficacy, and call for additional data across a wider range of soil textures, management practices, and hydrologic years, including wetter summers, before drawing strong conclusions on optimum P-application timing. \textcite{vanderSalm2012} argue that data spanning at least a normal, a wet, and a dry year are necessary to understand nutrient losses and transport routes. \textcite{SalcedaGonzalez2025} recommend more extended calibration and validation periods, longer treatment periods, and further data collection to obtain definitive evidence of total suspended solids reductions and more conclusive orthophosphate results under cover cropping. \textcite{Ulén2007} note that as long a period as possible is needed for trend estimation because water-quality changes may occur more slowly than the P-loss changes driven by agricultural measures. Long-term or continued monitoring is also recommended to evaluate P leaching to groundwater and nutrient removal of green-land infiltration systems \parencite{Duan2015}, to determine whether enhanced effluent-driven P losses persist \parencite{Monaghan2004}, and to confirm bioreactor and saturated-buffer effectiveness across different soils, landscapes, and weather years \parencite{Jaynes2019}.

Experimental-design refinements are proposed alongside longer records. \textcite{Schilling2013} recommend establishing a suitable correlation among paired basins across a range of climate conditions, restricting paired sampling to the March--July period when most water and nitrate flows occur, delaying post-BMP monitoring to enhance the likelihood of detecting change, and using BMPs capable of achieving at least the minimum detectable change. \textcite{Sunohara2016} argue that BMP impacts on downstream receptors should be assessed simultaneously at multiple scales, from field to watershed. \textcite{Irvine2019} call for within-field and within-plot studies to elucidate the mechanistic processes responsible for nutrient export at the sub-watershed scale, and for additional field-scale studies within high-export sub-watersheds. \textcite{Xin2016c} recommend further experiments under a wider range of conditions to test the applicability of derived effective-yield formulas and models, while \textcite{Xiong2025} call for refined $^{15}$N loss observation and larger trial areas more representative of actual production, and \textcite{Loper2024} see mesocosms as an opportunity to quantify nutrient-loss responses to management changes.

\subsection{Nitrogen management: rates, timing, tools, and loss pathways}

Optimising N application rates is expected to continue receiving primary research and regulatory focus, and accurate 4R implementation will require site-specific knowledge to balance agronomic and environmental goals \parencite{Christianson2015b}. \textcite{Bandura2025} propose concrete standards for the next generation of N-rate tools: they should be developed and refined for various application timings using the economically optimum N rate (EONR) as a benchmark, judged adequate if they predict N need within 10~kg~N~ha$^{-1}$ of EONR, and considered environmentally unacceptable if they exceed EONR at an N:corn price ratio of 5.6 by more than 30~kg~N~ha$^{-1}$. \textcite{Liu2021} call for more accurate N-fertilisation recommendations that consider both environmental and economic benefits, effective local nutrient-management plans, case-by-case cost analyses, and coordinated agency action (e.g.\ USDA-NRCS with USEPA) to address the economic losses from lower fertiliser rates and to promote 4R adoption. \textcite{Banger2018} advocate additional field experiments across wider crop-management practices and soil types, particularly to identify potential yield tradeoffs from N-rate reductions, and note that if such yield-loss risks appear, greater emphasis should shift to cover crops and edge-of-field practices (bioreactors, wetlands, buffers) to meet load-reduction targets. \textcite{Gupta2023} suggest precision fertiliser application based on soil N availability as a new approach to prevent N losses via subsurface drainage.

Timing and method of application are recurring open questions. \textcite{Liang2018} propose combining timing with rate to further optimise N management and note the value of models such as WHCNS for evaluating short- and long-term drainage, N-loss, and crop-growth effects. \textcite{Hertzberger2019} call for additional coordinated field research on nutrient timing, note the continuing evolution of the MANAGE database as a research resource, and emphasise the need for greater research on P transport in subsurface drainage. \textcite{Malone2007b} highlight the need for more research to confidently quantify differences between split and single spring N applications and to check possible over-prediction of winter nitrification and N loss for fall manure on soybean. \textcite{HernandezRamirez2011} note it remains unknown whether reduced-rate fall manure could match sidedress or preplant applications environmentally, and point to nitrification inhibitors, winter cover crops, more complex rotations, and multifunctional landscapes as additional mitigation avenues, with controlled drainage as a supplementary method. \textcite{Elmi2005} stress the need to quantify all major N loss pathways simultaneously and to obtain better information on the relative merits of manure versus mineral management systems.

Local-scale N testing and decision approaches also require refinement. Late-spring nitrate-test (LSNT) programmes are framed as an initial start to be replaced or refined as new research emerges, with high-clearance application equipment to reduce delayed-sidedress risk, and cropping diversity or edge-of-field denitrification needed to reach drinking-water goals \parencite{Jaynes2004}. \textcite{Malone2010} treat their conclusions as an impetus for further research, suggesting weather-based model refinement into simple tools and improved understanding of year-to-year soil mineral-N variability to reduce soil-testing needs and increase side-dressing adoption. Grain-quality and grain-protein monitoring is recommended where LSNT is implemented, along with risk insurance or other benefits for adopters \parencite{Karlen2005}. Broader gaps include comprehensive volatilisation-loss studies to determine loss extent and controlling factors, targeted research from data-poor regions, and provision of ancillary climate, yield, and soil data to support predictive modelling and policy \parencite{Chatterjee2020b}. \textcite{Haghnazari2020} call for research comparing economic and environmental consequences of N-fertiliser use under global-warming projections and for identifying best field-management practices to reduce environmental impacts.

\subsection{Phosphorus transport processes and mitigation research gaps}

Phosphorus in subsurface drainage is repeatedly identified as under-researched relative to nitrogen. \textcite{Hertzberger2019} explicitly call for greater research effort on P transport in subsurface agricultural drainage and better understanding of nutrient-related water-quality issues. \textcite{Yuan2022} note that more research is needed on conservation-practice effectiveness in reducing dissolved-P loss from agricultural fields, particularly in subsurface flow, and on the causes of annual variability in reduction rates. \textcite{Pease2018} identify increasing understanding of the pedologic factors influencing P transport in subsurface-drained landscapes, and the effect of management on P loss, as key research needs across the Eastern Corn Belt, requiring examination of individual and combined management factors. \textcite{deOliveira2022} recommend refining P-loss frameworks through additional within-field studies, continued data collection on soil testing, site history, agronomic management, and water-quality/flow monitoring, and better understanding of the "where" and "when" of drainage P loss to sharpen conservation-practice timing and placement for disproportionate, targeted benefit.

Several studies target specific dissolved and bioavailable P questions. \textcite{Uusitalo2003} argue that directing runoff into the subsurface pathway and establishing field-margin buffers are only first steps, that mitigating bioavailable particulate P requires reducing particle mobilisation across the entire field surface, and that controlling diffuse dissolved reactive P (DRP) is a longer-term task than reducing erosion. \textcite{Uusitalo2018} state that assessing the pace of P stratification requires greater data, suggest simulation tools such as the Annual P Loss Estimator (APLE) for evaluating mixing intervals, and note it remains an open question whether no-till adoption as a eutrophication-control measure is justified at their site. \textcite{Ni2020} call for quantifying dominant mechanisms of agricultural conservation practices (ACPs) and site-specific conditions under complex scenarios, more data on other ACPs affecting surface and subsurface DRP and total P, and studies of optimal soil-test-P levels under changing climate to address STP--precipitation interactions. \textcite{Wang2018} note further study is needed to confirm whether results apply to manure-injection scenarios and to test whether pre-application or immediate post-application tillage decreases DRP losses through tile drainage. Manure-related P risk is also flagged: the risk of immediate P loss after a single manure application needs fuller understanding \parencite{Madison2014}. Modelling refinements proposed include adjusting manure-P mineralisation parameters and winter drainage, setting manure decomposition rates as adjustable parameters, and further investigating tillage effects using daily or monthly datasets \parencite{Pan2023}, alongside continued efforts to improve P-loss modelling and to represent crop stress prior to maturity \parencite{Wang2019}.

Additional P-loss frontiers concern event dynamics, ratios, and legacy sources. \textcite{Brendel2019} recommend designing BMPs that reduce P losses during frequent small events (which can drive extreme P and suspended-solids losses), understanding relationships between flow, precipitation, transport pathway, and P fraction at the catchment scale, and using these to target BMP implementation. \textcite{Schilling2017b} call for refining organic-P/total-P concentration and load ratios for Iowa and the US Midwest, evaluating their sensitivity to load-estimation method and seasonality, quantifying wastewater-treatment-plant contributions, and identifying specific sources of variability such as manure applications and animal densities. \textcite{Jarvie2017} ask whether changing climate patterns contribute to increased soluble reactive P delivery, note it is unknown how much of a 40\% reduction target could be achieved by 4R practices alone, and warn that drawing down legacy P and disconnecting tile-flow P loads at times of maximum receiving-water sensitivity may be increasingly costly and intrusive. \textcite{Ulén2007} add that knowledge of the availability of eroded P to algae is still limited. Cost-effective P-mitigation hierarchies are proposed, prioritising decreasing soil Olsen P and targeting alum applications to critical source areas that lose the most bioavailable P from a small part of the farm \parencite{McDowell2015}, and P removal via regular grass harvesting under continuous-grass tile-drained fields \parencite{Zhang2015b}.

\subsection{Bundling, sequencing, and edge-of-field practices}

A pervasive conclusion is that no single BMP suffices and that combinations must be studied, sequenced, and optimised. \textcite{Sunohara2016} state plainly that no one BMP is a silver bullet and that multiple BMPs are likely required to complement controlled tile drainage. \textcite{Lenhart2016} argue that multiple BMPs in series are needed to achieve nutrient-reduction goals, that combinations of in-field and edge-of-field practices show great promise but have been little studied, and that wetlands and other water-retention BMPs must be developed to more effectively control eutrophication in sensitive downstream waters, with vegetation establishment encouraged before diverting water into a wetland. \textcite{Christianson2014} conclude that attaining water-quality goals will require a combination of practices implemented at a high level and in a targeted manner. \textcite{Makhtoumi2025} recommend combining conservation tillage with cover crops (and harrowing the top ${\sim}10$~cm to disturb leaching pathways), pairing in-field with edge-of-field conservation to maximise nutrient removal, and incorporating legacy effects to assess the real potential of practices. \textcite{Kreig2019} call for examining synergistic effects among conservation practices, formalising optimisation studies, and addressing farmer concerns about 2-tier versus 3-tier riparian buffers. \textcite{McLellan2015} emphasise better understanding of practice interaction, greater scientific understanding of the nitrogen removal capacity (NRC) of various practices, more data from stream-restoration projects over longer periods, and improved mapping of tile-drained and ditched land to evaluate available land area for N-removal practices.

Region-specific bundling agendas reinforce this. \textcite{Baulch2019} lay out a broad cold-region research programme: investigating drainage effects on nutrient export under diverse weather while comparing surface, tile, and undrained fields; studying controlled drainage, channel design, and maintenance through wet--dry cycles; assessing rotational and spring tillage with altered nutrient management in reduced-till systems; determining net effects of perennial-forage conversion (stand age, composition, snowmelt infiltration); quantifying nutrient retention and remobilisation in reservoirs, wetlands, ponds, and buffers; refining nutrient management by minimising late-fall and winter applications and targeting hydrologically connected hotspots; and developing tile-drainage mitigation (controlled drainage, bioreactors, physico-chemical amendments) suited to snowmelt hydrology. \textcite{Gupta2023} call for more robust, cost-effective BMPs to aid cover-crop benefits and for robust edge-of-field BMPs to mitigate climate warming. \textcite{Husk2024} suggest wider application of agricultural bioreactors, P reactors, controlled drainage, and municipal low-impact development, potentially alongside internal-loading management such as sediment dredging, P inactivation, biomanipulation, aeration, and nutrient precipitation. \textcite{HernandezRamirez2011} and \textcite{GAYNOR1995b} point to intercrops and management that enhances crop nutrient uptake and alters watershed hydrology as promising for P-transport reduction, with tillage options beyond ridge and zero tillage still to be investigated. Edge-of-field structure research is also called for: quantifying gains from individual components of saturated riparian buffers, installing monitoring wells and backflow valves, and clarifying suspended-solids sources \parencite{Singh2024b}, and confirming that saturated-buffer N removal increases as buffer vegetation matures \parencite{Jaynes2019}. \textcite{Spiese2023} call for novel BMPs to address both nitrate and DRP simultaneously.

\subsection{Spatial targeting, hot spots, and critical source areas}

Closely related is the frontier of targeting practices to the parts of the landscape that matter most. \textcite{Melland2008} recommend identifying hydrologically active parts of the landscape and managing them as priorities for mitigating surface-runoff P, avoiding higher stocking rates and fertiliser inputs in waterlogged runoff-generating areas. \textcite{Irvine2019} describe ongoing work to determine whether elevated loads originate from numerous low-intensity hot spots or a few key hot spots, seeking to identify land-management-specific hot spots and targeted BMPs. \textcite{Speir2022} find that targeted cover-crop implementation may be required to reduce soluble reactive P export, that regional saturation of cover crops on all fallow ground will likely be needed for consistent large annual reductions in nitrate-N and SRP, and recommend subwatershed-based optimisation. \textcite{Baulch2019} and \textcite{McDowell2015} likewise stress targeting critical source areas and hydrologically connected hotspots. \textcite{Brendel2019} note their results assist in targeting areas for BMP implementation and remediation. \textcite{Makhtoumi2025} emphasise that further progress in prioritisation promises to improve return on investment in conservation practices. \textcite{Kvítek2023} prioritise automating spatial-targeting principles into an easily applicable GIS tool, projecting results into river-basin management plans, and recommend whole-basin surveys with SWOT and cost-benefit analyses and stakeholder engagement.

\subsection{Cropping systems, cover crops, perennials, and land use change}

Improved cropping practices are widely proposed as under-explored levers. \textcite{Elli2023} recommend exploring genetic-trait impacts on water-table depth, yields, and soil-N dynamics; investigating how trait changes affect water tables and reduce climate-projection uncertainty; and examining additional strategies including irrigation, cultivar maturity groups, further cropping systems, off-season cover crops, perennial systems, green manure, and stacked multi-management strategies extended to further environments. Cover-crop optimisation is a recurring need: \textcite{deOliveira2022} call for optimising cereal-rye management (early planting for fall growth, targeting biomass accumulation) to assess DRP-loss reductions and minimise P-loss tradeoffs, while \textcite{PrezaFontes2023} suggest growing cover crops longer in spring to prevent increases in N$_2$O emissions and adjusting incentives to include enhanced-efficiency N sources. \textcite{McNeill2026} note that information on different winter-cover types is needed to inform BMPs for the increasingly hydrologically active non-growing season, that "win-win" BMPs benefiting both soil and water are needed, and that because over-winter BMP effects differed for N versus P, practices targeting individual nutrient loss pathways must be identified. \textcite{Krueger2013} recommend conservation practices beyond rye cover cropping because subsurface nitrate remained high and rye had no effect on DRP or drainage volume, suggesting French drains to replace surface inlets and applying part of crop nutrients via irrigation water. \textcite{Mitchell2000} note cover cropping requires innovative management before soil-quality benefits are expressed in productivity, along with better understanding of emergence impedance and stubble-reinforced crusts.

Perennial and land-use-change options feature prominently. \textcite{Huang2023} suggest that gaining additional soil organic carbon may require purposeful management, that poorly drained depressions may be good candidates for conversion to perennial bioenergy crops such as \textit{Miscanthus $\times$ giganteus}, or restoration to wetlands by obstructing tile drains and planting native species. \textcite{Norberg2022} conclude that mitigating N losses is most effectively achieved by keeping soil covered year-round with a growing crop such as perennial ley, while the focus for P should be on decreasing surface runoff. \textcite{Arrueta2025} call for additional data on alfalfa's effect on subsurface discharge during large events, across diverse management and geographies, and better understanding of how nutrient and transport factors couple to govern temporal variability in export. \textcite{Djiemon2019} argue for better cropping practices to ensure long-term subsoiling returns, integrating cover crops and autumn cereals to dry soils in spring, and controlling machinery traffic and weight. \textcite{Woodbury2018} call for research on the value and payments for other externalities (P and sediment reduction, greenhouse-gas mitigation, biodiversity), which may further highlight switchgrass benefits when replacing maize. \textcite{Siksnane2017} list a suite of measures warranting evaluation: fertilisation planning and nutrient balancing, converting arable land to grassland, intercultural cultivation, winter crop farming, controlled drainage, organic farming, buffer-zone creation, and artificial wetlands with phytoremediation. \textcite{GARDNER1992} recommend greater emphasis on soil types beyond red-brown earths, increased subsoil root growth, deep ripping and deep placement of nutrients and gypsum, crop species that open subsoils, and lucerne in dryland systems with rotational grazing. \textcite{Oquist2007} note more research across regions is needed to determine whether conversion to alternative farming practices improves water quality elsewhere.

\subsection{Water-table and drainage-water management}

The operation and siting of controlled drainage, DWM, and subirrigation constitute a further frontier. \textcite{Jaynes2012} are unaware of research optimising DWM operation for the soils, climate, and crops of central Iowa, and call for identifying locations where a greater field area falls under control-structure influence and for incentive programmes that share installation and management costs to increase adoption. \textcite{Adeuya2012} note more experience is required to determine optimum dormant-season management depth. \textcite{Cooke2012} require longer studies to determine DWM's effect on yield. \textcite{Gunn2015} recommend better on-site monitoring plans, long initial calibration periods, anti-seep collars on new and retrofitted systems to minimise seepage around control structures, and educational programmes to help farmers understand management-plan recommendations. \textcite{Gunn2019} call for extending subirrigation studies to more diverse regions with in-depth economic assessment, assessing variable subirrigation strategies across locations, and evaluating the importance of decreased yield uncertainty. \textcite{Baker2004} note more research is needed on whether soil-aggregate property differences affect the long-term sustainability of water-table management by subirrigation. Nutrient-trading programmes are identified as needed to increase adoption of shallow drains \parencite{Ghane2021b}, and \textcite{HernandezRamirez2011} note controlled drainage could in principle be implemented as a supplementary method for reducing drainage flow and nutrient losses. In paddy and rice systems, new water-management options may change the whole agronomy and require new infrastructure and better farm management \parencite{DarziNaftchali2014}, and integrated water-and-nutrient management systems with advancing precision agriculture are encouraged to support open-field smart agriculture \parencite{Jung2024}.

\subsection{Models, indices, tools, and data infrastructure}

Numerous studies frame the frontier in terms of improved models, indices, and shared data. \textcite{Buczko2007} outline development of a phosphorus index for Germany, starting regionally (e.g.\ Mecklenburg-Vorpommern), incorporating German conditions onto a Pennsylvania-style backbone, using factors such as wind erosion, subsurface drainage and soil leaching, groundwater transport in peat soils, and stream-network connectivity, and possibly combining the P index with an N index for concentrated animal operations. \textcite{Pärn2012} argue nutrient-index schemes should be improved by including wetland and riparian-buffer removal capacity, incorporating buffer N-removal efficiency or catchment soil wetness into both N and P index models. \textcite{Negm2016} note that regression-based tools for nutrient-credit trading via DWM can be developed for other US and international regions, that Midwest equations are underway, and that user-friendly web interfaces and educational material will be needed. \textcite{Wang2016} discuss using steeper stress s-curves in critical periods and milder transformations elsewhere, noting aeration-stress effects on yield may be more easily identified at plot scale against field-collected yield data. \textcite{Boivin2006} call for focus on hydraulic-conductivity and shrink-swell parameters and their seasonal time-dependence, direct in-field measurement methods, more comprehensive long-term datasets, and delineation of seasonal preferential-flow processes. \textcite{ROBINSON1992} note yield-loss coefficients can be adjusted with regional data and that their methodology is applicable to other regions with similar terrain. \textcite{Adeuya2012} note their statistical analysis could apply to similar time-series data. Data-resource development is emphasised through the evolving MANAGE database \parencite{Hertzberger2019}, provision of raw and ancillary data \parencite{Chatterjee2020b}, and coupling models with tools such as INVEST to assess synergies and tradeoffs among water quality and other ecosystem services \parencite{McLellan2015}.

\subsection{Economics, incentives, policy, adoption, and stakeholder engagement}

A consistent frontier is that biophysical understanding alone is insufficient; economic evaluation, incentive design, policy reform, and farmer engagement are decisive for adoption. \textcite{Liu2021} call for advanced case-by-case cost analyses to guide policymakers, farm owners, and researchers, and for inter-agency action to promote 4R implementation. \textcite{Yuan2022} note more data are needed on the cost-effectiveness and ecological benefits of ACPs. \textcite{McLellan2015} plan to incorporate practice costs and cost-effectiveness into future work. \textcite{Christianson2014} recommend continued positive incentives such as CREP funding for wetlands (given high public and lower perceived private benefit), more information and educational programmes to increase adoption (especially for crop rotation), and greater consideration of broader-scale, non-excludable benefits in multifunctional landscapes. \textcite{David2015} argue that incentivising on-farm research and collaboration among farmers, programmes that generate a collective sense of conservation need, and scaling individual efforts to meso- and macro-scale policy are essential. \textcite{Jaynes2012} and \textcite{Karlen2005} both call for cost-share and risk-insurance mechanisms to raise adoption. \textcite{Kvítek2023} call for reassessment of legislation and revised incentives for the 2024+ period to remove barriers to natural-water-retention-measure adoption. \textcite{Marja2013} recommend national policy change to retain open ditch-and-field networks where subsurface drainage prevails, noting that removing state support for ditch replacement may suffice, with targeted agri-environmental support as an additional incentive in high-revenue regions. \textcite{Ghane2021b} identify nutrient-trading programmes as needed to increase shallow-drain adoption.

Adoption, capacity, and social dimensions are stressed across contexts. \textcite{Okwany2016} call for detailed cost-benefit analyses, technical assessment of the most cost-effective designs, economic and social studies of upscaling potential, joint farmer--researcher experimentation, and greater information dissemination and capacity building, noting adoption depends on farmers' long-term cost-benefit perspective. \textcite{Ghumman2011} emphasise creating know-how and farmer confidence before government project execution, with farmer participation playing a very important role. \textcite{Ward2018} identify key practice-use issues: where different BMPs should be applied, availability of tools to assist producers, how best to bundle BMPs, and transferability of research knowledge to commercial farms, and note farmers need additional education on soil ecology and biology. \textcite{Cuevas2019} stress that proposed measures must consider stakeholders' willingness to adopt. \textcite{Saikkonen2014} argue urgently for research integrating ecological, pollution-related, and economic impacts of alternative land uses to develop environmentally sustainable yet economically profitable agriculture, and for models addressing optimal policy design. In salinity-affected and irrigation contexts, \textcite{Kotb2000} propose a package spanning reduced and consolidated rice plantings, increased irrigation staff, soil-fertility monitoring, farmer compensation, flexible cost-recovery for drainage projects, public-awareness programmes, and gradual subsidy liberalisation. \textcite{Sellars2022} conclude that improved feasibility of drainage-water-recycling requires markedly higher, less volatile irrigated yields and lower adoption costs, that producers may prefer alternative adaptations such as drought-resistant seeds or precision technology, and that internalising large welfare benefits would positively influence adoption.

\subsection{Greenhouse gases, emerging contaminants, and multi-objective tradeoffs}

An expanding frontier integrates greenhouse-gas (GHG) and contaminant considerations into drainage-mitigation assessment. \textcite{PrezaFontes2023} call for evaluating net GHG emissions in longer-term experiments. \textcite{Carstensen2021} argue that the fate of waterborne methane leaving via the riparian zone warrants further research and that GHGs should be included alongside phosphorus balances when comparing and selecting mitigation measures. \textcite{Burton2024} note more study is needed to understand the magnitude and fate of nitrous-oxide losses. \textcite{HernandezRamirez2011} call for comparative life-cycle assessment of agroecosystems based on complete N balances encompassing all ecosystem services and environmental impacts. Emerging contaminants and biological interactions are also flagged: \textcite{Spiese2023} call for better constraining the mechanism and lifetime of glyphosate in soil and water and long-term DRP datasets across land-use types to elucidate glyphosate's effect. \textcite{Schaafsma2019} argue that planter modifications must be investigated to reduce exposure of aquatic non-targets to neonicotinoid and other seed-applied pesticides. \textcite{Knorr2023} call for nutrient-mitigation strategies readily implemented without adding time or cost burdens on producers and effective across seasons, and for expanded research on how watershed microbial communities respond to and control nutrient fluxes. \textcite{Fraindová2022} call for studies of water-quality changes during different rainfall-runoff events in relation to landscape characteristics and anthropogenic activities to identify compound sources and protect water quality.

\subsection{Salinity, reclamation, and irrigation-return management}

In arid and salinity-affected drainage systems, distinct research and management needs emerge. \textcite{Sharma2004} suggest sequential reuse across cereals, grasses, forest plantations, and halophytes for future research, and innovative, cost-effective evaporation-pond designs for effective containment of salt and toxic elements. \textcite{MANN1982} caution that only small amounts of gypsum should be applied during early reclamation to avoid excessive losses, while \textcite{He2015} recommend chemical amendments (gypsum, elemental sulphur, possibly agricultural lime) in long-term management of problem soils and consideration of more environmental factors for field assessment. \textcite{Cuevas2019} foresee increased use of tolerant rootstocks and varieties, screening of halophyte germplasm for recovering salinised lands, and transgenic crops as proposed solutions. \textcite{Christen2010} propose combining land-treatment systems with pond treatment, restricting hydraulic loading to allow soil drying, storing excess wastewater for off-season application, using larger land areas with longer irrigation intervals, providing better pre- or post-treatment to reduce biochemical oxygen demand, growing and harvesting potassium-accumulating herbages, and reducing sodium in the waste stream. \textcite{Duan2015} call for pre-treatment (dilution or nutrient removal) before wastewater application and for controls to minimise N leaching by enhancing plant uptake and preventing preferential flow. \textcite{Mitchell2000} note detailed economic evaluations of integrated reuse systems are needed. \textcite{Purkey2001} raise questions about the ultimate use and suitability of retired land overlying a shallow water table surrounded by irrigated fields. \textcite{Bjorneberg2020} prioritise continued reduction of sediment loss (especially from furrow-irrigated fields), improved sediment removal in water-quality ponds, improved irrigation management to reduce return flow, and irrigation and nutrient management to limit nitrogen leaching. \textcite{Tian2023b} recommend selecting biochar with good amelioration, low contaminant content, and low environmental risk, and urgently assessing its long-term positive and negative impacts on soil quality and crop yield. \textcite{Zhang2015} call for determining optimum rates of organic soil conditioners to minimise adverse water-quality impacts while improving soil quality and productivity.

\subsection{Emerging pressures, downstream endpoints, and remaining infrastructure questions}

Finally, several studies point to under-studied downstream endpoints, climate-driven hydrologic shifts, and physical drainage-technology questions. \textcite{Apsite2017} note that projected shifts of maximum discharge from spring to winter must be considered in future flood planning, that discharge changes may alter seasonal water supply for hydropower, and that further effects will occur. \textcite{Husk2024} note further long-term nutrient drawdown and internal-loading management may be needed to achieve additional trophic improvement in receiving lakes. \textcite{Winikoff2020} call for identifying and testing novel management practices to control hybrid cattail and reed canary grass while protecting native perennials, with targeted control needed to maintain excavation benefits. \textcite{PALKO1994} suggest mass-balance studies for further insight and direct liming of ecologically important watercourses. \textcite{Lucci2014} propose developing Bayesian networks to explore trade-offs of different wintering strategies and note stock access to waterways as a complex P source outside current scope. On drainage-material and infrastructure design, \textcite{Zhang2025b} plan to explore relationships between perforation characteristics, drainage performance, and compressive performance, noting further field validation is needed, while \textcite{ASGHAR1995} call for comparing permeameter results across Pakistan, Belgium, and Canada, standardising the number of drain-plate perforations, and investigating the effect of perforation size. \textcite{Zhang2022b} note that future field experiments should include N loss among essential observations. \textcite{Eastman2010} report ongoing efforts to demonstrate total-P transport mechanisms under different BMP scenarios. \textcite{Grenon2023} call for further long-term comparative field studies with more hydrologic and climatic data to identify the long-term environmental effects of water-management systems on P, alongside simultaneous implementation of improved fertiliser management, precision agriculture, and vegetated buffer strips. \textcite{Hashemi2020} recommend special management strategies in heavy-textured subsurface-drained soils, adoption of water-table and fertilisation management for sustainability, and assessment of new technologies before introduction to paddy fields. Collectively, these frontiers underscore that advancing tile-drainage water-quality management will require sustained, multi-scale, and interdisciplinary research linking process understanding, targeted and bundled practices, robust tools, and economically and socially viable pathways to adoption.